\section*{Acknowledgment}\addcontentsline{toc}{section}{Acknowledgment}

\subsection*{Individual contributions}

This is a single-student project. I, Anish Talla, chose the research
question, designed the comparison, selected the dataset and the metric,
wrote the project specification that every later decision was checked
against, made every research decision, ran and checked every experiment,
verified the numbers and citations in this report, and wrote and approved
the final text. No other student took part. My instructor,
[INSTRUCTOR NAME] of [INSTITUTION], [describe the instructor's role in one
or two sentences: for example, reviewed the study design before the test
split was opened, read drafts of the report, and advised on presentation.
Delete this bracket once filled in].

\subsection*{Use of artificial intelligence tools}

The Award asks for a detailed and truthful account of AI use. Two kinds of
AI appear in this project and they should not be confused. GPT-4o is the
object of study: it is the zero-shot model evaluated as System C and
System D, and every call to it is logged in the raw record described
below. Separately, AI coding and writing assistants were used as tools
while doing the work. That second use is what this section discloses.

\textbf{Tools and versions.} Anthropic Claude Code, the command-line
coding agent, was used from 28 July 2026, when the repository was created,
through 13 September 2026. The underlying models over that period were
Anthropic's Claude 4.x series and, in September 2026, Claude Fable 5.1
[confirm the exact model names from the session history before
submitting]. OpenAI Codex was used once, in late August 2026, on the
five-page conference version of this work. GPT-4o (\texttt{openai/gpt-4o}
through the OpenRouter API, default temperature) is System C and System D
and was not used as a writing or coding tool.

\textbf{Stages and purposes.} Claude Code was used at the following stages.
(1) Code development throughout: it drafted, under my direction and
review, the data-loading and evaluation code, the object-identity
reconstruction, the three systems, the training scripts including the
second-round recipe with cosine decay, weight averaging, and test-time
augmentation, the System D candidate generator, renderer, prompt, and
runner, and the statistical scripts for the object-clustered bootstrap and
permutation tests. I specified what each script had to do, read the code,
and checked its outputs against independent hand calculations; the three
bugs described in the report were caught this way. (2) Object-boundary
labeling: in reconstructing object identities, an automated first pass by
Claude assigned a decision and a confidence to each of the 251 ambiguous
frame boundaries. In a blind 30-boundary audit its calls agreed with mine
on 86.7\% of confident cases and 40.0\% of uncertain ones, so I decided
every uncertain call and every high-risk ``different object'' call
myself; this is reported in the Methods. (3) Analysis and diagnosis: it
helped audit the first draft's statistics, which is how the misdefined
angle-error figure and the single-seed weakness of the first training
round were found and corrected, and it implemented the second training
round and System D under a design I set and pre-registered. (4) Writing:
it helped edit, shorten, and restructure prose that I had drafted, and it
drafted some passages of the second-round and System D sections from the
verified results, all of which I revised and approved. OpenAI Codex was
used only to help restructure and shorten the conference version from my
existing report, code, and results.

\textbf{Timing and frequency.} Use was frequent: most working sessions
between 28 July and 13 September 2026 involved the coding assistant, in
sessions that ranged from minutes to several hours. The heaviest use was
in the last week of July and first week of August 2026 (building the
pipeline and the three systems) and on 12 and 13 September 2026 (the
second training round, System D, and this report's revisions).

\textbf{What AI did not do.} No AI tool chose the research question,
the dataset, the metric, or the hypotheses. No AI tool wrote the main
body of this report from nothing; the text was drafted by me and edited
with assistance. No experimental result was generated, altered, or
fabricated by an AI tool: every number in this report comes from scripts
in the repository whose outputs are saved, and the raw GPT-4o replies for
Systems C and D are kept verbatim in a log of 2,415 calls. No literature
was cited on the strength of an AI suggestion alone; I read each cited
paper and checked each reference.

\textbf{Verification materials.} Under ``Other Materials'' I am
submitting the session transcripts of the Claude Code conversations
exported from the tool, the raw GPT-4o call log, the source code
repository, and the saved outputs behind every table. [Export the
transcripts from \texttt{\textasciitilde/.claude/projects/} before
submitting and confirm this sentence matches what was uploaded.]

\subsection*{Other assistance and resources}

The Cornell Grasping Dataset was used under its public release. Model
weights for ResNet and Faster R-CNN came from the torchvision library.
GPT-4o calls were paid for through an OpenRouter account at a total cost
of about ten US dollars. Training ran on a personal Apple M4 laptop.
Experiment tracking used Weights \& Biases. No funding, laboratory,
commercial training center, or paid service other than the API and
tracking accounts named above contributed to this work.
